\section{Comparison of algebraic sheets and extension of terminal time}
\label{sec:comparison}

We prove a comparison principle for two small lists of scalar functions
and apply it to extend the terminal charts of an exact problem in time.
The comparison forces sheets that agree on many sampled incidences to
agree on a whole neighborhood.  Algebraic continuation then carries
these local comparisons to a common complex point, where finite jets
provide labels for representative equations.  This is how a fit on a
longer time block can remain known from the old terminal observation.

The comparison uses fresh conditional samplings.  Successful paths are
kept as unnormalized submeasures; the sampling kernels are never
conditioned on a later step succeeding.

We use real quantifier elimination and semialgebraic cell decomposition
with parameters in fixed dimension and degree: formulas of fixed format
define semialgebraic sets, and the numbers of connected components and
isolated points of their fibers have bounds uniform in their real
coefficients.  A one-variable semialgebraic function which is finite and
bounded on compact subintervals of \([2,\infty)\) has at most polynomial
growth at infinity.  These foundational forms, including their
parameter-uniform versions, may be found in
\citet[Chapters~3, 5, 7, and 14]{BasuPollackRoy}.

\subsection{Uniform norms and continuation for algebraic sheets}

A holomorphic function will be called an \emph{algebraic sheet of degree at
most \(D\)} if it satisfies
\[
 P(z,f(z))=0
\]
for a nonzero polynomial \(P\) of total degree at most \(D\).
There is no lower bound on a nonzero coefficient of \(P\).
All balls and polydiscs in the next lemma have fixed geometry.  In particular,
the real normalization ball has nonempty real interior and stays a fixed
distance inside the complex domain.

The next lemma belongs to the Bernstein--Remez theory of algebraic
functions; compare \citet[Sections~3.3--3.5]{RoytwarfYomdin1997}.
We include a proof of the formulation used here.

\begin{lemma}[Algebraic analytic norm comparison]
\label{lem:algebraic-norm}
Let \(\Omega\subset\C^n\) be a fixed ball or polydisc, and let
\(B\subset\Omega\cap\R^n\) be a fixed closed ball with nonempty interior.
Fix a compact set \(K\Subset\Omega\).  For algebraic sheets of degree at most
\(D\), normalized by \(\|f\|_B=1\), the following bounds are uniform:
\begin{enumerate}
\item for every fixed integer \(q\),
\(\max_{|\alpha|\le q}\|\partial^\alpha f\|_K\le C_q\);
\item at every \(z\in K\),
\[
 \max_{|\alpha|\le D}|\partial^\alpha f(z)|\ge c_K>0.
\]
\end{enumerate}
Consequently, let \(B_1\subset\Omega\cap\R^n\) be another fixed
closed real ball with nonempty interior.  For every measurable
\(E\subset B\) with \(\theta=|E|/|B|>0\),
\begin{equation}
 \|f\|_{B_1}
 +\max_{|\alpha|\le q}\|\partial^\alpha f\|_{B_1}
 \le C_q\theta^{-C}\|f\|_E,
 \label{a02:remez}
\end{equation}
The statements are homogeneous without normalization.  They are unchanged,
with the corresponding scaled derivatives, by affine translation and
rescaling of the base, including translation to a parallel real slice
through a complex point.  A sufficiently long finite jet at the center of
an interior ball controls the norm on that ball.
\end{lemma}

\begin{proof}
We first prove local boundedness of the normalized family, allowing the
degree in the function variable to drop in a limit.  Suppose \(f_\nu\) is
such a family.  Normalize the coefficient vector of a relation
\(P_\nu(z,W)\) to have norm one and pass to a subsequence on which it
converges to a nonzero polynomial \(P\).  This limiting polynomial must
involve \(W\).  Indeed, if \(P=P(z)\), then for every \(x\in B\), boundedness
of \(f_\nu(x)\) gives
\[
 P(x)=\lim_{\nu\to\infty}P_\nu(x,f_\nu(x))=0.
\]
A polynomial vanishing on the real interior of \(B\) is zero, a
contradiction.

Write \(a(z)\) for the leading nonzero coefficient of \(P\) as a polynomial
in \(W\), and set \(Z=\{a=0\}\).  On a compact subset of \(\Omega\setminus Z\),
all roots of \(P(z,\cdot)\) lie in a fixed disk.  A sufficiently large
circle in the \(W\)-plane is therefore root-free, uniformly over that
compact set.  The same circle remains root-free for \(P_\nu\) for large
\(\nu\).  Although some roots of \(P_\nu\) may escape to infinity, a
continuous chosen root cannot cross this circle.

Choose a point in the interior of \(B\setminus Z\).  The chosen value
\(f_\nu\) there is bounded.  It follows, by continuation along a compact
path and a neighborhood of that path in \(\Omega\setminus Z\), that
\(f_\nu\) is uniformly bounded there.  The complement
\(\Omega\setminus Z\) is path connected: one can join two points by
complex line segments and make small detours around the finitely many
zeros met on generic complex lines.  The preceding bound therefore
holds locally throughout that complement.

It remains to cross \(Z\).  At a point of \(Z\), choose a complex line
on which \(a\) is not identically zero, and a small disk in that line
whose boundary avoids \(Z\).  After translating the disk slightly, all
of its boundary circles still avoid \(Z\).  These circles form a compact
subset of \(\Omega\setminus Z\), where the bounds have already been
proved.  The maximum principle on the disks now bounds \(f_\nu\) in a
neighborhood of the original point.  A finite cover gives local uniform
boundedness on every compact subset of \(\Omega\).

Cauchy's estimates and compactness yield a subsequence converging
holomorphically on compacta, together with all fixed derivatives, to a
function \(f\).  Uniform convergence on \(B\) preserves \(\|f\|_B=1\);
in particular \(f\) is not zero.  The relation \(P(z,f(z))=0\) persists.
Divide \(P\) by its largest polynomial factor which is a power of \(W\).
The resulting relation has a nonzero constant coefficient \(a_0(z)\):
\[
 a_0(z)=-\sum_{k\ge1}a_k(z)f(z)^k.
\]
If \(f\) vanishes to order \(r\) at a point, then \(a_0\) vanishes to
order at least \(r\) there.  Since \(a_0\) is a nonzero polynomial of
degree at most \(D\), we have \(r\le D\).  Compactness, also for a sequence
of possible violating points in \(K\), proves both asserted uniform
bounds.

For completeness, these jet bounds imply a quantitative sublevel
estimate.  At each point of a fixed real interior ball, some directional
derivative of order at most \(D\) of either \(\Re f\) or \(\Im f\) has
absolute value at least a fixed positive constant.  One may use a fixed
finite set of real directions, by finite-dimensional norm equivalence
for symmetric tensors.  Bounds on one additional derivative make this
lower bound persist on a neighborhood of fixed radius.  If the order
is zero, that neighborhood contains no sufficiently small sublevel
values.  If the order is \(p\ge1\), restrict to segments parallel to
the selected direction.  From a subset of length \(m\) choose \(p+1\)
points separated by at least a constant times \(m\).  The divided
difference formula and the mean value theorem imply
\[
 |\{s:|f(s)|\le u\}|\le C u^{1/p}.
\]
Integrating over transverse slices and using a finite cover gives,
for some fixed \(c>0\),
\begin{equation}
 |\{x\in B_1:|f(x)|\le u\|f\|_{B_1}\}|
 \le C u^c |B_1|,\qquad 0<u<1.
 \label{a02:algebraic-sublevel}
\end{equation}
The same compactness argument compares the norms on the fixed
interior balls involved here.  Apply \Cref{a02:algebraic-sublevel}
with \(B\) in place of \(B_1\) to \(E\subset B\).  This bounds
\(\|f\|_B\) by \(C\theta^{-C}\|f\|_E\); norm comparison and
the derivative bounds then give \Cref{a02:remez} on \(B_1\).
Finally, the uniform lower bound on the
finite jet, after normalization on the comparison ball, proves the
last assertion.
\end{proof}

We record closure properties needed when using this lemma.  Sums,
differences, and fixed-degree polynomial evaluations of a fixed number
of algebraic sheets have polynomial relations of bounded degree.
Over the rational function field in \(z\), form the product of the
linear factors corresponding to all combinations of roots and clear
denominators.  Symmetry in each collection of roots bounds the degrees
of the resulting coefficients.  Restrictions to affine complex lines
also have bounded-degree relations.  If direct specialization makes
a relation identically zero, perturb the line on an interior disk,
normalize the coefficients of nontrivial specialized relations, and
take a limit.  Derivatives of the simple linear or quadratic sheets
used below have this property as well: implicit differentiation gives
rational expressions in the base variables and the same radicals,
with nonzero denominators on the comparison domain.  Eliminating the
radicals and clearing these denominators gives the required relations.
Thus \Cref{lem:algebraic-norm} applies to all the differences,
evaluations, and fixed derivatives used here.

Local comparisons will eventually be recorded at a common complex
point.  The next lemma controls the cost of moving a sheet along a
path that avoids its branch points, uniformly even when the equation
coefficients degenerate.

\begin{lemma}[Polynomial cost of continuation]
\label{a02:algebraic-path}
Fix the base dimension, polynomial degrees, number of sheets, and a
bound on the number of pieces of a broken straight path.  Work in a
fixed bounded complex base region.  The sheets may be affine functions,
roots of linear equations, or simple roots of quadratic equations in
the function variable whose quadratic coefficient is constant.
Differences and fixed-degree polynomial evaluations of these sheets
are allowed.

Suppose the linear coefficients needed for linear roots and the
discriminants needed for quadratic roots are nonzero throughout an
\(A^{-1}\)-neighborhood of the path, where \(A\ge2\).  On endpoint
comparison balls of radius \(c/A\), with fixed margins inside that
neighborhood, the norms of a continued function compare in both
directions with loss \(A^C\).  Fixed derivatives are controlled by
these norms, and sufficiently many endpoint jets control the norms,
with losses of the same form.  The exponent \(C\) is independent of
the values of the equation coefficients.
\end{lemma}

\begin{proof}
There are two separate uniformity arguments.  For a fixed \(A\), cover
the path by \(O(A)\) overlapping balls of radius comparable to \(A^{-1}\),
with fixed enlargements in the prescribed neighborhood.  Each branch
is simple and holomorphic on these enlargements.  Its polynomial
evaluations have bounded-degree relations, by the preceding closure
properties.  Applying \Cref{lem:algebraic-norm} on successive overlapping
balls gives a finite bound \(C^{O(A)}\) for the endpoint norm ratio.
It is uniform in all coefficient values and is bounded for \(A\) in
compact subintervals of \([2,\infty)\).  A function with zero starting
norm vanishes under continuation, so it need not enter a ratio.

We next show that the best such bound has a fixed semialgebraic
description.  Split all complex parameters into real and imaginary
parts.  The equation coefficients, the boundedly many path vertices,
and the endpoint branch choices form a fixed-dimensional parameter
space.  Nonvanishing on the tube is expressed by real quantification
over a segment parameter and a displacement of size less than \(A^{-1}\).

For a square root of a nonvanishing discriminant, continuation of the
endpoint signs can be encoded with finitely many crossings.  Choose a
generic ray from zero, avoiding the discriminant values at the path
vertices.  Along each straight piece the discriminant is a polynomial
in a real segment parameter.  Except when its argument is constant,
there are only boundedly many critical directions: differentiate the
argument, whose numerator is a real polynomial of bounded degree.
Choose the ray away from those directions; a constant-argument piece
can also be avoided.  Crossings of the ray are then transversal and
their number is bounded by the degree.  Relative to a square root of
the ray direction, label an endpoint square root by the sign of its
imaginary part.  Each crossing changes this sign, and the parity of
the bounded list of crossings specifies exactly the continued
endpoint root.  Crossing parameters, their order, transversality,
and the parity condition can all be described by polynomial
conditions and real quantifiers of fixed complexity.

To evaluate a norm on an endpoint ball, append one straight segment
from its center to a variable point of that ball and use the same
description of continuation.  Quantification over that point and its
selected roots describes the norm inequalities.  It follows from
quantifier elimination that the supremum of the endpoint ratio over
all admissible data is a one-variable semialgebraic function of \(A\).
The first argument proves that this supremum is finite and locally
bounded.  The polynomial growth theorem for one-variable
semialgebraic functions therefore bounds it by \(A^C\); these are the
fixed-format semialgebraic facts recalled above.
The \(O(A)\) comparison balls were used only to prove finiteness and
are not part of this semialgebraic description.

Reversing the path proves the reverse comparison.  Applying
\Cref{lem:algebraic-norm} on the endpoint balls, and accounting for
their radius in derivatives, proves the jet assertions.
\end{proof}

In particular, a tiny constant quadratic coefficient causes no loss
outside the stated bound.  For example, the roots of
\(\varepsilon W^2+zW+1\) can have very different sizes.  When its two
branch points are closer than the prescribed tube radius, an
admissible path cannot separate them and change the relevant
continuation parity.  In general the coefficient-independent
fixed-\(A\) bound and the fixed-format continuation description
perform precisely this separation of a bounded chosen branch from
unused escaping roots.

\subsection{Two-color matching}
\label{a02:two-color-subsection}

The next lemma compares close values with separated derivative data.
At value accuracy \(N^{-1}\), two lists of size at most
\(N^{d+o(1)}\), with a fixed \(d<1\), cannot exhibit both
behaviors on a sampled incidence law with the stated marginal bounds.
The functions need not be algebraic; that restriction enters only
when we pass from derivative data to graph norms.

\begin{lemma}[Two-color matching]
\label{lem:two-color-matching}
Fix \(D\ge1\), an integer \(J>D\), and \(d_0<1\).
There is no sequence \(N\to\infty\) with the following properties.
There are two indexed lists of real \(C^{J+1}\) functions on a unit
interval, each of cardinality at most \(N^{d+o(1)}\), where \(d\le d_0\),
and all derivatives through order \(J+1\) are bounded by
\(K=N^{o(1)}\).  There is a probability law on events \((i_1,i_2,t)\)
such that, for each color \(\kappa=1,2\), its \((i_\kappa,t)\)-marginal
is bounded by
\begin{equation}
 K\,\pi_{\kappa,i_\kappa}\,dt,
 \qquad
 \sum_i\pi_{\kappa,i}=1,
 \label{a02:index-marginal}
\end{equation}
and every retained event satisfies
\begin{align}
 |f_{1,i_1}(t)-f_{2,i_2}(t)|&\le K N^{-1},\\
 \max_{1\le j\le D}
 |f_{1,i_1}^{(j)}(t)-f_{2,i_2}^{(j)}(t)|&\ge K^{-1}.
 \label{a02:matching-gap}
\end{align}
The impossibility is uniform with sufficiently small fixed power
losses in place of all displayed subpower losses, for these fixed
orders and \(d\le d_0\).
\end{lemma}

\begin{proof}
Suppose such a sequence exists.  We record the functions' jets at
successive time and value scales.  The list size bounds the number
of finest records.  On most nearby scales, an alternating path along
functions from the two lists has a polynomial endpoint map with a
nonzero Jacobian.  Its image forces more records than this bound
allows.  The resampling argument below controls the law of these
paths while retaining their full dependence on the preceding choices.
This use of alternating paths and an endpoint Jacobian is analogous
to the method of refinements for incidence operators; compare
\citet[Sections~3--4]{ChristQuasiextremals}.

All probability restrictions in the proof have subpower mass, until
the bounded-length switching experiment is introduced.  Normalizing
such a restriction changes \Cref{a02:index-marginal} only by a
subpower factor.  We enlarge \(K\) as necessary.

\paragraph{Jet scales and nested states.}
Logarithmically pigeonhole the derivative gaps through order \(J\).
After a subsequence their exponents are \(\beta_j\), meaning that
the \(j\)-th gap has size \(N^{-\beta_j+o(1)}\) on the retained events.
Exponents above a fixed large constant greater than one are cut off;
at a cutoff we use only the corresponding upper bound.  We have
\[
 \beta_j\ge0,\qquad \beta_0\ge1,\qquad
 \beta_j=0\ \text{for some }1\le j\le D.
\]
At value scale \(N^{-a}\), choose a time radius \(N^{-r(a)}\)
on which every Taylor term of the difference is at most \(N^{-a}\),
up to subpower factors.
The \(j\)-th term requires \(r(a)\ge(a-\beta_j)/j\).
The corresponding accuracy for an individual \(j\)-th derivative
is \(N^{-(a-jr(a))}\).  This exponent need not increase with
\(a\), so we retain its largest preceding value.  Thus, for
\(0\le a\le1\), define
\begin{equation}
 r(a)=\max_{1\le j\le J}\frac{a-\beta_j}{j},
 \qquad
 l_j(a)=\max_{0\le x\le a}\{x-jr(x)\}.
 \label{a02:jet-scales}
\end{equation}
Then \(r(0)=0\), and \(r\) is increasing, convex, and piecewise linear.
Moreover,
\begin{align}
 l_0(a)&=a,&
 0\le l_j(a)&\le\beta_j,&
 l_j(a)-l_{j+i}(a)&\le i r(a).
 \label{a02:translation-exponents}
\end{align}
The last inequality follows by evaluating \(l_{j+i}\) at a point
where the maximum for \(l_j\) is attained.  Since
\(r(x)\ge x/D\), for \(a>0\) we also have
\begin{equation}
 (J+1-j)r(a)>l_j(a).
 \label{a02:remainder-gap}
\end{equation}
Indeed, if a maximizing point \(x\) for \(l_j(a)\) is positive, use
\((J+1)r(x)>x\) and monotonicity of \(r\); if \(x=0\), then
\(l_j(a)=0\) and \(r(a)>0\).

On an open piece where the unique active term in \(r\) is \(m\),
convexity gives
\begin{align}
 r(a)&=(a-\beta_m)/m,\\
 l_j(a)&=a-jr(a)\quad(0\le j\le m),\\
 l_j(a)&>a-jr(a)\quad(j>m).
 \label{a02:active-piece}
\end{align}
For \(j\le m\), the function \(x-jr(x)\) is nondecreasing up to
the active piece.  For \(j>m\), moving slightly to the left within
that piece strictly increases it.  The active \(\beta_m\) has not
been cut off, because the cutoff terms cannot maximize \(r\) on
\([0,1]\).

We choose the finite grid before constructing its states or selecting
their atom masses.  Fix an integer \(C>5J\).  Trim a small fixed amount
from the ends of every active piece, also staying away from zero, and
omit pieces shorter than the trim.  On the remaining compact intervals,
the gaps
\[
 l_{m+i}(a)+i r(a)-\beta_m>0\quad(1\le i\le J-m)
\]
and all needed remainder gaps are uniformly positive.  Choose a common
small step \(b\), after these trims, and subdivide each remaining
interval.  Write \(A_0\) for its first subdivision point.  We will use
steps \([a,a+b]\) only when \(a+Cb\) still lies in that interval.
Take a finite increasing grid containing \(1\), all these subdivision
points, and every required \(a+b\) and \(a+Cb\).  The construction
that follows is performed anew for each such grid.  All its selections
and limits in \(N\) precede any decrease of \(b\) or of the trims.

For a level \(a\) in this fixed grid, the state \(F_a\) records the
dyadic interval \(I_a\) containing
the event time, of length comparable to \(N^{-r(a)}\), the bins of
both arms' derivatives of orders \(0,\ldots,J\) at the center of
\(I_a\), with widths \(N^{-l_j(a)}\), and all earlier grid states.
The dyadic partitions are nested.  Taylor translation over \(I_a\)
has uncertainty at derivative order \(j\) bounded by
\begin{equation}
 K\sum_{i=0}^{J-j}N^{-l_{j+i}(a)-ir(a)}
       +K N^{-(J+1-j)r(a)}
 \le K N^{-l_j(a)}.
 \label{a02:translated-bin-width}
\end{equation}
Here \Cref{a02:translation-exponents,a02:remainder-gap} handle,
respectively, the binned terms and the remainder.  Applying the same
calculation to the two arms, starting from the event gaps and
\(l_j\le\beta_j\), shows that their center bins differ by only \(K\)
neighbors, at every recorded level.

An index from one color and \(I_1\) determine that arm's entire
history.  The other arm's history then has only subpowerly many
neighbor choices.  Therefore
\begin{equation}
 |\supp F_1|\le N^{d+r(1)+o(1)}.
 \label{a02:state-support}
\end{equation}
For this fixed grid, make a further subpower selection on which the
masses of occurring atoms, measured \emph{before this selection},
are \(N^{-u_a+o(1)}\), uniformly at each grid level.  Extremely light
atoms can first be omitted using \Cref{a02:state-support}.  The
exponents are nonnegative and nondecreasing, and
\begin{equation}
 u_1\le d+r(1).
 \label{a02:state-upper}
\end{equation}
If a retained coarse atom at \(a\) has retained descendants at
\(c>a\), their number is at most
\begin{equation}
 N^{u_c-u_a+o(1)}.
 \label{a02:descendant-count}
\end{equation}
To see this without any assertion about later conditional
probabilities, denote the pre-selection law by \(p_{\rm pre}\).
Each retained child has \(p_{\rm pre}\)-mass at least
\(N^{-u_c-o(1)}\), and the disjoint children lie in a parent of mass
at most \(N^{-u_a+o(1)}\).  Summing their pre-selection masses proves
\Cref{a02:descendant-count}.  Write \(\Pp\) for the final normalized
event law.

The upper bound \Cref{a02:state-upper} will contradict the growth
forced by matching.  More precisely, on almost all small steps in
an active \(m\)-piece we will prove
\[
 u_{a+b}-u_a\ge b\left(1+\frac1m-o_b(1)\right).
\]
Summing these increments, then decreasing \(b\) and removing the
trims, will force the uniform upper bound \(d+r(1)\) to be at least
\(1+r(1)\).
The next two constructions establish this increment on most steps
of the already fixed grid.

\paragraph{A palette with little entropy growth.}
We use the entropy chain rule and conditional mutual information
in their standard forms; see \citet[Chapter~2]{CoverThomas2006}.
Fix \(D_1>5\).  The palette \(Z\) is the pair of \(m\)-th derivative
bins at the center of \(I_{A_0}\), now at width
\(N^{-\beta_m-D_1b}\).  Since \(l_m(A_0)=\beta_m\),
\[
 H(Z\mid F_{A_0})\le2D_1b\log N+O(1).
\]
At a step retained in the grid construction, let
\[
 L_a=\log\frac{\Pp(Z\mid F_{a+Cb})}{\Pp(Z\mid F_a)}
\]
at the sampled palette value.  Conditional mutual information and
the entropy chain rule give
\[
 \sum_a \E L_a\le C H(Z\mid F_{A_0}).
\]
The negative part has bounded expectation.  Indeed, for probability
vectors \(p,q\),
\[
 \sum_{p(z)<q(z)}p(z)\log\frac{q(z)}{p(z)}
 \le \frac1e\sum_zq(z)\le\frac1e.
\]
Apply this conditional on the fine state, with \(p\) its palette law
and \(q\) the coarse palette law.  There are order \(b^{-1}\) steps,
so the average of \(\E(L_a)_+\) is
\(O(b^2\log N+1)\).  Markov's inequality gives an average probability
\(O(b^{3/4})+o_N(1)\) for
\(L_a>b^{5/4}\log N\).  A second application shows that, except for
a proportion \(O(b^{1/2})+o_N(1)\) of the steps, this failure
probability is at most \(b^{1/4}\).  Call these steps good.  Thus at
a good step,
\begin{equation}
 \Pp(Z=z\mid F_{a+Cb})
 \le N^{b^{5/4}}\Pp(Z=z\mid F_a)
 \label{a02:palette-likelihood}
\end{equation}
at the sampled value, apart from probability tending to zero with
\(b\).

\paragraph{The switching experiment.}
Fix a good step.  We will make \(m+1\) traversals, whose \(m\)
internal arrival times supply the variables of a full-rank endpoint
map.  Its endpoint must lie in the small union of targets supplied
by \Cref{a02:descendant-count}; comparing those volumes will force
the required increment.  Start with \(E^0\sim\Pp\), keep every time
in the same \(I_a\), and proceed as follows.  At the beginning of
traversal \(p+1\), draw a fresh event \(\widetilde E^p\) from
\(\Pp\) conditional only on
\[
 F_{a+Cb}(\widetilde E^p)=F_{a+Cb}(E^p).
\]
Write \(t_p=t(E^p)\) for the running time and
\(\widetilde t_p=t(\widetilde E^p)\) for the refreshed event's
witness time.  Both lie in the same fine interval
\(I_{a+Cb}\), but the refresh leaves the running time at \(t_p\).
Choose a color \(\kappa_p\) by the rule below and set
\[
 \varphi_p=f_{\kappa_p,i_{\kappa_p}(\widetilde E^p)}.
\]
Draw \(E^{p+1}\), including its time \(t_{p+1}\), freshly from
\(\Pp\) conditional only on this arm's color, index, and \(I_a\).
Traverse \(\varphi_p\) from \(t_p\) to \(t_{p+1}\).  The increments
can have either sign.

Before any failures are removed, \(E^p\) and \(\widetilde E^p\)
have marginals bounded by \(2^p\Pp\).  Refreshing on a state preserves
a marginal bound of this form.  For a traversal, dominate the
adaptively selected arm kernel by the sum of the two color kernels;
each color kernel preserves \(\Pp\) after averaging over its
index and interval.  This proves the assertion inductively.

The possible coarse states stay in a subpower neighbor set of the
start.  A refresh preserves \(F_a\), since the finer state includes
it.  A traversal preserves the chosen arm's bins at \(a\) and at
all earlier levels: its index and \(I_a\) determine these data.
At either endpoint the other arm is within \(K\)-neighbor bins.
There are only finitely many coordinates at the fixed grid and
only \(m+1\) traversals.  Thus all possible coarse states lie in a
start-determined set of size \(K\), after increasing \(K\).

Let \(t_a\) be the center of \(I_a\), and subtract the order \(J\)
Taylor polynomial \(T\) of one starting arm there.  On traversal
\(p+1\), put \(\varphi=\varphi_p\) and set
\[
 s=(t-t_a)N^{r(a)},\qquad
 Y_j=N^{l_j(a)}\partial_t^j(\varphi-T)(t),
 \quad 0\le j<m.
\]
Along a traversal,
\begin{equation}
 \frac{dY_j}{ds}=Y_{j+1}\quad(j<m-1),\qquad
 \frac{dY_{m-1}}{ds}=U+O(N^{-2b}),
 \label{a02:controlled-jets}
\end{equation}
where \(U\) is a constant of size at most \(K\).
Indeed the factor in the last derivative is \(N^{\beta_m}\).
Taylor expansion of this residual \(m\)-th derivative at \(t_a\)
shows that its variation is bounded by
\[
 K\sum_{i=1}^{J-m}
 N^{\beta_m-l_{m+i}(a)-ir(a)}
 +K N^{\beta_m-(J+1-m)r(a)}.
\]
The strict trimmed gaps make this smaller than \(N^{-2b}\), for
\(b\) sufficiently small.

Choose \(U\) from the fresh palette, the higher jet bins at \(A_0\)
recorded in the coarse history, and the fixed frame \(T\).  Taylor
translation from the center of \(I_{A_0}\) to \(t_a\) gives an error
\(O(N^{-D_1b})\) in its \(m\)-th term, while the higher-bin errors
have exponents
\[
 l_{m+i}(A_0)+ir(A_0)-\beta_m>0.
\]
The remainder has the same strict gap as before.  This discretizes
\(U\) to the precision required in \Cref{a02:controlled-jets} using
only the start, the coarse state, the palette, and the color.
The two controls available in a fresh pair differ by at least
\(K^{-1}\): the \(m\)-th derivative gap at its event has exponent
\(\beta_m\), and all the approximation errors are smaller by a
fixed power.  Choose either color on the first leg.  Thereafter
at least one of the two available controls is separated from the
preceding control by \(1/(2K)\); choose such a color and enlarge \(K\).

At a switch, the old and new arms have \(K\)-neighbor fine bins
at \(a+Cb\).  The running time \(t_p\) and the witness time
\(\widetilde t_p\) lie in that finer interval, so the bins control
the jump at the actual switching time \(t_p\).  Applying
\Cref{a02:translated-bin-width} there gives a jump in \(Y_j\) of
size at most
\[
 K N^{l_j(a)-l_j(a+Cb)}
 =K N^{-(1-j/m)Cb}
 \le N^{-2b},\qquad j<m.
\]
Here \(C>5J\) leaves room for subpower factors.  This also accounts
for the comparison to the starting arm at the first switch.

Kill a step if \Cref{a02:palette-likelihood} fails at its refresh.
Also kill it if its chosen color-index/interval pair has
\(\Pp\)-mass less than
\(\pi_{\kappa,i}|I_a|/\log N\), with the weights from
\Cref{a02:index-marginal}.  For each color the sum of these latter
thresholds over all indices and intervals is \(1/\log N\).
The marginal domination just proved therefore makes the total
failure probability small.  On a retained pair, the new-time
conditional density is bounded by
\[
 \frac{K\pi_{\kappa,i}}
 {\pi_{\kappa,i}|I_a|/\log N}
 \le \frac K{|I_a|}.
\]
Consequently the \(m\) internal arrival times can also be required
to be pairwise separated by \(|I_a|/K\), with negligible additional
loss: test each new internal time against the preceding ones and
choose the final reciprocal-subpower separation threshold smaller
than the reciprocal of the preceding density bound.  The experiment
has successful probability bounded below when averaged over starts,
first for a sufficiently small fixed \(b\) and then for large \(N\).

\paragraph{Domination conditional on the entire past.}
Fix the exact start \(E^0\).  For every coarse state in its
start-determined neighbor set, push its conditional palette law
to each of the two possible discretized controls, using that
state's earlier data and the fixed frame \(T\).  Sum these
pushforwards to obtain a measure \(\Lambda\) on controls, with
\(\Lambda(\R)\le K\).

Let \(\mathcal H_p\) contain the entire sampling history before
the next refresh, including all hidden event data.  This history
fixes the current finer and coarse states.  The next refresh has
exactly the finer-state conditional law, independently of the
rest of \(\mathcal H_p\).  On its good palette values,
\Cref{a02:palette-likelihood} dominates that law by the corresponding
coarse palette law.  The chosen arm can depend on the full history
and the full refreshed event, but its point mass is bounded by the
sum of the two candidate-control point masses.  Finally, conditional
on that complete refreshed event and its actual chosen index, the
retained new-time density is at most \(K\,ds\).  Therefore the
\emph{unnormalized} successful step kernel satisfies
\begin{equation}
 Q_{\mathcal H_p}(dU\,ds)
 \le K N^{b^{5/4}}\Lambda(dU)\,ds
 \quad\text{for every admissible }\mathcal H_p.
 \label{a02:full-past-kernel}
\end{equation}
The same \(\Lambda\) works for every such history with this start.

For a nonnegative necessary test depending only on the start and
the successive controls and times, \Cref{a02:full-past-kernel} may
be iterated backwards.  Integrate the final step with its full
previous history fixed.  The upper integral then depends only
on the previous displayed controls and times, so repeat with the
previous step.  No successful kernel is renormalized and no kernel
is conditioned on future survival.  This is why possible hidden
dependence of future events on the refreshed line identity does
not change the product upper bound.

\paragraph{The endpoint map and its target.}
Write \(U_1,\ldots,U_{m+1}\) and \(s_1,\ldots,s_{m+1}\) for the
controls and arrival times, put \(s=s_{m+1}\), and keep the starting
time \(s_0\) fixed.  Solve \Cref{a02:controlled-jets} with exact
constant controls and no switching errors.  Its endpoint differs
from the actual low jets by \(O(N^{-2b})\).  At fixed final time
\(s\), the control contribution to its \(j\)-th coordinate is
\begin{equation}
 \frac1{(m-j)!}
 \left[
 U_1(s-s_0)^{m-j}
 +\sum_{p=1}^m(U_{p+1}-U_p)(s-s_p)^{m-j}
 \right].
 \label{a02:endpoint-polynomial}
\end{equation}
The omitted part depends only on the initial jets and \(s\).
The formula follows by telescoping the integrals of constant
controls; it remains valid for signed increments.

For fixed controls the map from
\((s_1,\ldots,s_m,s)\) to the final time and low jets is polynomial
of fixed degree.  Differentiating \Cref{a02:endpoint-polynomial}
with respect to \(s_p\) gives
\[
 -\frac{U_{p+1}-U_p}{(m-j-1)!}(s-s_p)^{m-j-1}.
\]
Thus the absolute Jacobian determinant is exactly
\begin{equation}
 \frac{\displaystyle
 \prod_{p=1}^m|U_{p+1}-U_p|
 \prod_{1\le p<q\le m}|s_p-s_q|}
 {\displaystyle\prod_{k=0}^{m-1}k!}.
 \label{a02:endpoint-jacobian}
\end{equation}
It is at least \(K^{-1}\) on successful paths.  The number of
preimages on this regular locus is bounded in terms of the degree
and dimension: each is an isolated solution of a fixed-degree
polynomial system, and the fixed-format semialgebraic fiber bound
applies uniformly in the coefficients.

Set \(c=a+b\).  A possible final child state \(F_c\) records its
own interval \(I_c\), and the final time lies in that interval.
At each such time, translate each final-arm jet bin from the
center of \(I_c\), subtract \(T\), and apply the coarse normalization.
The uncertainty in derivative order \(j<m\) is bounded by
\begin{align}
 &K N^{l_j(a)}
 \left(
 \sum_{k=j}^J N^{-l_k(c)-(k-j)r(c)}
 +N^{-(J+1-j)r(c)}
 \right) \notag\\
 &\hspace{35mm}\le
 K N^{l_j(a)-l_j(c)}
 =K N^{-(1-j/m)b}.
 \label{a02:final-target-widths}
\end{align}
This uses the child's interval, not the whole coarse interval.
It includes every higher-jet uncertainty and the Taylor remainder.
The simulation error \(N^{-2b}\) fits these widths.  The normalized
final-time interval has length \(O(N^{-b/m})\).  Hence, for either
choice of final color, the target associated with one child has
volume at most
\begin{equation}
 K N^{-b\left(\frac1m+\sum_{j=0}^{m-1}(1-j/m)\right)}
 =
 K N^{-b\left(\frac1m+\frac{m+1}{2}\right)}.
 \label{a02:target-volume}
\end{equation}
Its center may move rapidly with the final time; Fubini's theorem
uses only the widths of the fibers in \Cref{a02:final-target-widths}.

Take the union of these targets over both colors and all retained
children of every possible coarse neighbor of the start.  By
\Cref{a02:descendant-count}, it has at most
\(N^{u_{a+b}-u_a+o(1)}\) constituent targets.  Membership in this
union, together with the internal-time and consecutive-control
separations, is a necessary test depending only on the start,
controls, and times.  Other hidden success conditions can be
discarded for the upper bound.

For each fixed separated control tuple, the change-of-variables
formula, \Cref{a02:endpoint-jacobian}, and the bounded regular-fiber
multiplicity bound the time integral by \(K\) times the target
volume.  This estimate is uniform in the controls.  We may therefore
integrate it against \(\Lambda^{m+1}\), even when \(\Lambda\) is
atomic, and use \Cref{a02:full-past-kernel}.  The conditional
successful probability is at most
\[
 N^{u_{a+b}-u_a
 -b\left(\frac1m+\frac{m+1}{2}\right)
 +(m+1)b^{5/4}+o_N(1)}.
\]
The averaged lower bound for success gives, on every good step,
\begin{equation}
 u_{a+b}-u_a
 \ge b\left(1+\frac1m-o_b(1)\right),
 \label{a02:entropy-increment}
\end{equation}
since \((m+1)/2\ge1\).

For this fixed trimmed grid, first let \(N\to\infty\), passing to
a subsequence that fixes the set of good steps.  Sum
\Cref{a02:entropy-increment}; all omitted increments are nonnegative.
Only after this fixed-grid limit do we let \(b\to0\), so that the
proportion of excluded steps tends to zero.  Then let the trims tend
to zero.  The sum of the
\(b/m\) terms tends to
\(\int_0^1r'(a)\,da=r(1)\), and the sum of the \(b\) terms tends
to one.  Thus \Cref{a02:state-upper} implies \(d\ge1\), a
contradiction.  The homogeneous exponents can depend on the fixed
grid; only their uniform upper bound is used when passing to
finer grids.

Finally, if no uniform small fixed-power tolerance existed, choose
counterexamples with loss exponents tending to zero and \(N\)
tending to infinity.  Their lists can be bounded using \(d_0\),
and they would form the subpower counterexample sequence just
excluded.  This proves the last assertion.
\end{proof}


\begin{corollary}[Local matching of scalar graphs]
\label{a02:local-graph-matching}
Fix the base dimension, algebraic degrees, real comparison balls
with fixed holomorphic margins, and \(d_0<1\).  There is
\(\eps_*>0\) with the following property.  Let \(M,\Delta>0\), and put
\(R=M/\Delta\).
For sufficiently large \(R\), two lists of real-on-real algebraic
sheets cannot satisfy all of the following, with
\(0<\eps\le\eps_*\) and \(d\le d_0\):
\begin{enumerate}
\item each list has at most \(R^{d+\eps}\) members, and each member
is within \(M R^\eps\) in norm of one common algebraic reference
sheet;
\item a subprobability law on pairs of indices and a common real
base point in the comparison domain has total mass at least
\(R^{-\eps}\), with each index-base marginal bounded by
\(R^\eps\pi_{\kappa,i}\,dz\), where the \(\pi_{\kappa,i}\) are
probability weights, separately in each color;
\item on its events, the two values differ by at most
\(\Delta R^\eps\), whereas the norm of their difference on the
comparison ball is at least \(M R^{-\eps}\).
\end{enumerate}
Fixed multiplicative constants may be absorbed by decreasing
\(\eps_*\) or taking \(R\) larger.
\end{corollary}

\begin{proof}
The positive pair mass makes both lists nonempty.  Choose a member
\(f_*\) of one list as reference.  By the common-reference bound
and the triangle inequality, every list member satisfies
\(\|f_i-f_*\|\le2MR^\eps\).  Subtract \(f_*\) and divide by \(M\);
this reference has the same algebraic degree bound and common
holomorphic domain as the list members.
The closure properties and \Cref{lem:algebraic-norm} give derivative
upper bounds through any required fixed order with loss \(R^{C\eps}\).
At every incident point, the difference has a jet of bounded order
of size at least \(R^{-C\eps}\).  The zeroth-order value cannot
supply this lower bound when \(R\) is large and \(\eps\) is
sufficiently small, because its upper bound is \(R^{-1+\eps}\).
A finite set of real directions detects a positive-order
directional derivative.  Select one order, direction, and interior
product box carrying a fixed fraction of the pair mass.

Slice parallel to that direction.  Fubini's theorem supplies a
fiber interval with sliced pair mass at least \(R^{-C\eps}\).
Before normalizing this slice, each of its two index-time
marginals is still bounded by \(R^{C\eps}\) times the original
probability index weights.  Normalization costs only another
\(R^{C\eps}\).  Rescale the fiber interval, whose geometry and
holomorphic margins are fixed.  The restricted sheets meet
\Cref{lem:two-color-matching} with \(N=R\), with only fixed
multiples of \(\eps\) in all loss exponents.  Choose \(\eps_*\)
below that lemma's tolerance divided by these constants.
\end{proof}


\subsection{Trading terminal thickness for time}
\label{a02:terminal-extension-subsection}

We now apply scalar-sheet matching to a true terminal chart system.
The aim is to enlarge its time intervals by a factor \(N_0^e\),
while losing at most \(Ce\) in thickness depth. The new chart labels
must still be known from the old terminal observation. This last
requirement determines the construction: one old incidence will choose
a new label, and a second incidence on the same trajectory will provide
a short list containing it.

\begin{lemma}[Terminal time extension]
\label{lem:terminal-extension}
Fix \(d_0<1\) and angular parameters \(S,\eta>0\).
There are constants \(C<\infty\) and \(c>0\), depending only on
these fixed parameters and the fixed model dimensions and test
degrees, with the following property. Suppose an exact problem in
\(C(d)\), \(d\le d_0\), of length \(L\) has a true terminal system
at horizon exponent \(h>0\). For every
\[
 0<e<\min(h,cL)
\]
there is, on a retained subsequence, a known atlas at horizon
exponent \(h-e\), whose labels are known at the old observation
\(p_L\), and whose fitting thickness is \(a_{L'}\), where
\[
 0<L'<L,\qquad L-L'\le Ce.
\]
The resulting charts, trajectory assignments, tests, and selected
incidences are tempered in the original exact sequence. The
constant \(C\) is independent of the particular exact problem,
its length, and its tempered complexity exponent.
\end{lemma}

We recall the quadratic identities used to replace implicit tests by
graphs. Write \(z\) for the base variables of a test \(P(z,w)\)
of total degree at most two, and put \(c=P_{ww}\), which is constant.
Then
\[
 \mathcal D=P_w^2-2cP
\]
is independent of \(w\). If \(P,P_w,c\) are real, \(P_w\ne0\), and
\(2|cP|\le\frac12|P_w|^2\), then \(\mathcal D>0\), and the
real root \(f\) for which \(P_w(z,f)\) has the sign of
\(P_w(z,w)\) satisfies
\begin{equation}
 |w-f|
 =\frac{2|P(w)|}{|P_w(w)|+\sqrt{\mathcal D}}
 \le\frac{2|P(w)|}{|P_w(w)|}.
 \label{a02:nearest-quadratic-root}
\end{equation}
For \(c=0\), use the corresponding linear formula. For \(c\ne0\),
the affine critical center \(F_{\rm ctr}\) satisfies
\begin{equation}
 w-F_{\rm ctr}(z)=\frac{P_w(z,w)}{P_{ww}}.
 \label{a02:critical-center}
\end{equation}
These formulas follow by completing the square. A nonvanishing
discriminant on a simply connected complex ball supplies a simple
holomorphic branch there; a linear equation with nonzero coefficient
supplies its branch by division.

\begin{proof}
Throughout this proof \(N=N_0\), so all subpower factors refer to
the original exact sequence. Put \(\delta=N^{-L}\). In each world
replace \(w\) by \(Bw\); the derivative unit is then one, and
terminal hidden bins have width \(\delta\), up to dyadic rounding.
We undo this change at the end. Every constant implicit in
\(N^{O(e)}\) is independent of the exact problem and of \(L,h,e\).
For fixed \(e\), subpower factors can be absorbed into such losses.
The construction constants will be fixed before the final restriction
on \(e/L\).

The proof has three parts. We first compare the old equations at
two incidences in a longer time block. We then record their agreement
by finitely many derivatives at a common complex point and choose
one real representative equation for each derivative bin. Finally,
we fix one source time per block, use the source incidence to assign
trajectories to these bins, and normalize the representative equations
to obtain the new charts.

\par\smallskip\noindent\textbf{1. Comparing the old terminal equations.}

Prepare the true terminal system by \Cref{a01:saturation}, so its
labels are known by \(p_L\) and its true floors persist. Let
\(q_0=N^{-h+o(1)}\) be its dyadic time length. Select whole labels
and homogenize their spatial singular values. By largest-axis
fullness these are, up to subpower factors, \(q_0,q_0g\) in
version two, where \(g=N^{-\beta}\) and \(\beta\ge0\). In version
one put \(g=1\) and use the single spatial radius \(q_0\). Set
\begin{equation}
 n_Q=n(L)q_0^j g.
 \label{a02:terminal-floor-scale}
\end{equation}
Each used old label has selected incidence mass at least
\(n_Q/K_0\) in normalized \(q_0\)-block time. We use these
floors to count old labels; subsequent restrictions of paired
incidences need not retain a floor for each label.

Choose a coarser dyadic partition into blocks \(I\) of length
\(q_1\), with \(q_1/q_0\) comparable to \(N^e\). This is possible
because \(e<h\). In a world \(\gamma\) and block \(I\), sample
a trajectory from its original prior \(\nu_\gamma\), and sample
\(t_0,t\) independently and uniformly in \(I\). Retain the pair
when both times are selected incidences of the old terminal system.
Weight this experiment by \(\omega_\gamma q_1\).

Let \(s_I(\ell)\) be the selected time fraction of trajectory
\(\ell\) in \(I\). The total pair mass is
\[
 \sum_{\gamma,I}\omega_\gamma q_1
       \int s_I(\ell)^2\,d\nu_\gamma(\ell).
\]
By Jensen's inequality it is at least the square of the old
selected mass, and hence is bounded below by an inverse-subpower
factor. Since \(0\le s_I\le1\), either one-time marginal of
this unnormalized law is dominated by the old selected incidence law.

Write \(A\) and \(A'\) for the old labels at \(t_0\) and \(t\),
including their \(q_0\)-subblocks. Their tests, in the normalized
hidden coordinate, are \(P_A\) and \(P_{A'}\). We call \(A\) the
source label. Its own incidence supplies
\[
 |P_{A,w}(t_0,y(t_0),w(t_0))|\ge K_0^{-1}.
\]
The label \(A'\) is known from the observation at \(t\); its own
derivative floor is at \(t\).

Here is the form of the labels we will construct. From \(A\) we
will obtain a coarse moving spatial box \(D\) and a simple graph
near the source point. For each \(D\), a common complex point
\(\alpha_D\) will serve to compare these graphs: after continuation,
we bin a fixed finite tuple of their derivatives at \(\alpha_D\).
Call this bin \(b\). The pair \(c_0=(D,b)\) will determine a
representative real equation, independently of the old source
identity. On the retained paired incidences, \(A'\) will determine
a subpower list of possibilities for \(c_0\). Only a logarithmic
derivative-size refinement will later be added to the output label.

At this stage \(t_0\) remains averaged. In particular the estimates
used to choose a representative equation will hold before one fixes
the source time. Once that time and the complex comparison points
are fixed in a block, the source construction will assign at most
one \(c_0\) to each trajectory throughout the block.

In one world and new block the old label count is at most
\begin{equation}
 \frac{K_0N^e}{n_Q}.
 \label{a02:old-label-count}
\end{equation}
Indeed the new block contains \(O(N^e)\) old subblocks, and within
each the old floors and disjoint assignments give at most
\(K_0/n_Q\) labels. All pair masses below within this world and
block use \(\nu_\gamma\) and the two normalized uniform time
priors; restrictions are left unnormalized.

Polynomial extrapolation from an old subblock to the new block
gives, for either old test on its assigned trajectory,
\begin{equation}
 \sup_I|P|\le\delta N^{C_0e},
 \qquad
 \sup_I|P_w|\le N^{C_0e},
 \label{a02:old-test-extrapolation}
\end{equation}
after increasing a fixed \(C_0\).  The polynomials on a line have
degrees at most two and one, respectively, in both model versions.
Spatial deviations from the old moving chart center also enlarge
by at most \(K_0 O(N^e)\).

For comparison at the source, we also need to evaluate the witness
test away from its own incidence.  The derivative of each test is
at least \(N^{-e}\) in absolute value at the opposite sampled time,
except on power-small absolute pair mass.  To prove this, fix the
trajectory and its
own selected incidence first.  The derivative is affine in time
and has maximum at least \(K_0^{-1}\) on \(I\).  An affine
function of maximum \(M\) is below \(\varepsilon M\) in absolute
value on a subset of relative length \(O(\varepsilon)\).
The other time is still drawn from the uniform prior; imposing
its selection can only reduce this bad mass.  Taking
\(\varepsilon=K_0N^{-e}\) and treating both colors proves the claim.

Cap the one-time densities at both vertices on a fixed
thickness-depth grid of spacing at most \(e\), including \(0,L\),
by
\[
 N^e n(r).
\]
Also cap the joint terminal/angular density at angular depth
\(u=Re\) by
\[
 n(L)N^{e-SRe},
\]
where \(R\) is a fixed constant to be chosen shortly and
\(Re\le\eta L\).  The typical bounds and marginal domination
make these deletions power-small.  They are analytical masks
used for estimates of unnormalized measures.

\paragraph{A common spatial frame.}
In version two, consider a pair whose old narrow normals differ
in sine angle by more than \(gN^{Re}\); there is something to
prove only when \(gN^{Re}<1\).
The velocity lies in two strips of widths \(K_0g\), because the
old spatial chart bounds hold throughout their respective
subblocks.  Their intersection has diameter at most
\(K_0N^{-Re}\).

We bound the exceptional pair mass using the angular cap before
any conditioning on that exception.  Fix \(A\).  At the opposite
time, the spatial point lies in its extrapolated rectangle of
area at most
\[
 q_0^2gN^{C_0e}.
\]
At each exact base, apply \Cref{lem:quadratic-sublevel} with
the value width from \Cref{a02:old-test-extrapolation}, the
opposite-time derivative floor \(N^{-e}\), and hidden-bin
width \(\delta\).  This gives at most \(N^{O(e)}\) terminal
hidden bins.  Enlarge \(C_0\), independently of \(R\), so that
the count is at most \(N^{C_0e}\).  Given the terminal observation, the old
knowledge of \(A'\) leaves only \(K_0\) choices.  Each disagreeing
pair confines the velocity to \(K_0\) angular bins at depth
\(Re\).  For the opposite-time marginal of the unnormalized pair
law, integrate the capped density over the spatial rectangle and
these hidden and angular bins, then sum over \(A\).  The bound is
\[
 \frac{K_0N^e}{n_Q}
 \bigl(q_0^2gN^{C_0e}\bigr)
 \bigl(N^{C_0e}\bigr)
 \bigl(n(L)N^{(1-SR)e}\bigr).
\]
The four factors are, respectively, the old-label count, spatial
area, hidden-bin count, and joint density cap; the partner-label
and angular-bin multiplicities are absorbed in \(K_0\).
The normalized time integral is \(\int_I dt/q_1=1\), and
\(n_Q=n(L)q_0^2g\) in version two.  Thus the exceptional mass is at most
\begin{equation}
 K_0N^{(2+2C_0-SR)e}.
 \label{a02:normal-exception}
\end{equation}
Choose \(R\) large enough that this is power-small.  In version
one no normal comparison is needed.

Put \(g^\#=\min(1,gN^{Re})\), with \(g^\#=1\) in version one.
We now choose a coarse moving box \(D\), deterministically from
\(A\), of spatial radii \(q_1,q_1g^\#\) in version two, and
\(q_1\) in version one.  When \(g^\#<1\), grid the unoriented
normal at spacing \(g^\#\).  In that rounded frame grid the
normal component of the old center velocity at spacing \(g^\#\),
and grid its position at the new block midpoint at spacings
\(q_1,q_1g^\#\).  Use the resulting midpoint and normal velocity,
and take zero tangential center velocity.  When \(g^\#=1\), use
fixed axes and a static midpoint grid of spacing \(q_1\).
The entire spatial trace of the line on \(I\) lies in this
coarse moving box up to subpower factors.

Moreover, \(A'\) determines a list of only \(K_0\) possible
\(D\)'s on retained pairs.  Enumerate the boundedly many possible
rounded orientations of \(A\) neighboring the normal of \(A'\).
For each, project the midpoint and center velocity of \(A'\)
in that same rounded frame and enumerate their subpower
neighborhoods in the stated grids.  To justify these
neighborhoods, compare both old centers to the common line.
Their midpoint errors have norm at most \(K_0q_1\) and old
normal components at most \(K_0q_1g\); their velocity errors
are at most \(K_0\) and \(K_0g\), respectively.  A rotation
error of size \(g^\#\) acts on these differences, giving
\(K_0q_1g^\#\) and \(K_0g^\#\).  It does not act on an
uncontrolled absolute midpoint.  This proves the stated
compatibility lists.

Use coordinates \(z=(\tau,\xi)\) for \(D\): rescale time in \(I\)
to a unit interval, subtract its moving center, and divide the
spatial coordinates by the displayed radii.  A common subpower
dilation makes the relevant traces lie in a fixed box, with
bounded speeds in normalized time.  The spatial Jacobian back
to physical coordinates is at most
\begin{equation}
 N^{C_1e}q_0^jg.
 \label{a02:coarse-jacobian}
\end{equation}
Here \(C_1\) can depend on the already fixed \(R\).
Write \(z_0\) for the base point at \(t_0\).

Let \(J_D\) be the spatial Jacobian of these coordinates, including
the common dilation.  For a fixed entry \((A,D)\), or
\((A',D)\), push the restricted unnormalized pair law to
\(z_0=(\tau_0,\xi_0)\), integrating out the other sampled time.
Its density with respect to \(d\tau_0\,d\xi_0\) is bounded by
\begin{equation}
 N^{C_1e}n_Q.
 \label{a02:source-base-density}
\end{equation}
Indeed, the source's own test or the witness's extrapolated test
at \(t_0\), together with the respective derivative lower bound,
allows at most \(N^{C_0e}\) hidden terminal bins by
\Cref{lem:quadratic-sublevel}, with the same value width and
derivative floor as in the preceding count.  The capped
one-time density is \(N^en(L)\), and
the change of base variables satisfies
\[
 \frac{dt_0\,dy}{q_1}
   =J_D\,d\tau_0\,d\xi_0.
\]
Consequently the density is at most
\(J_DN^en(L)N^{C_0e}\le N^{C_1e}n_Q\), after increasing
\(C_1\) in \Cref{a02:coarse-jacobian}.  No conditioning on the
label and no new lower mass bound for that entry is used.

We will repeatedly use the following counting consequence.
Suppose a group construction gives at most \(N^{C_1e}\)
options per old label.  Omitting entries of mass less than
\[
 n_QN^{-C'e}
\]
costs a power-small total amount if \(C'\) is a sufficiently
large fixed constant.  This follows by multiplying the
threshold by the option count and \Cref{a02:old-label-count},
in each world/block, and then summing with the original
weights.  Apply this first to the entries \((A,D)\) and
\((A',D)\).  We retain their pre-deletion masses when using
the resulting lower bounds.

\paragraph{Stable simple graphs and comparison mass.}
We now construct the graphs whose agreement will later be recorded
at the anchor.  Write \(P_{A,D}\) for the source equation expressed
in \(D\)'s real affine base coordinates, after the initial
normalization \(w\mapsto Bw\), and use \(P_{A',D}\) for the
witness equation in the same coordinates.  These equations have
total degree at most two and constant second \(w\)-derivative.
For either equation write \(P\) temporarily, and choose an early
constant \(C_A\)
larger than the value and derivative losses in
\Cref{a02:old-test-extrapolation}.  If
\[
 |P_{ww}|\ge\delta^{-1}N^{-C_Ae},
\]
replace this test for comparison by
\(\widetilde P=w-F_{\rm ctr}(z)\).  \Cref{a02:critical-center} and the derivative upper bound
give
\[
 |w-F_{\rm ctr}(z)|\le\delta N^{O(e)}
\]
throughout the new block.  Otherwise let \(\widetilde P=P\).
Denote the resulting equations by \(\widetilde P_{A,D}\) and
\(\widetilde P_{A',D}\).  The conversion, including the curvature
test and any affine replacement, depends only on the old label and
\(D\), before a source time or a trajectory is fixed.
For every unreplaced test, the choice of \(C_A\) makes
\(|P_{ww}P|\) much smaller than \(|P_w|^2\) at the tested
point \(z_0,w(t_0)\), where \(|P_w|\ge N^{-e}\).
Thus
\begin{equation}
 N^{-O(e)}\le
 \mathcal D(z_0):=
 \widetilde P_w^2-2\widetilde P_{ww}\widetilde P
 \le N^{O(e)}.
 \label{a02:source-discriminant}
\end{equation}
For replaced tests the discriminant is one.

For each converted equation, the coefficients of its discriminant
polynomial in \(z\) are bounded by
\(N^{O(e)}\), uniformly over the retained entries.
Indeed, the entry's pre-deletion mass is at least
\(n_QN^{-C'e}\), whereas its source-base density is at most
\(n_QN^{C_1e}\).  Its projection therefore has Lebesgue measure
at least \(N^{-(C'+C_1)e}\).  On that projection the upper bound in
\Cref{a02:source-discriminant} holds.  Polynomial Remez
therefore bounds the coefficients and hence the discriminant
and its derivatives on a fixed complex enlargement.
This argument uses the original heavy entry, before later
joint deletions.

Fix a net of real balls \(U\) of radius \(N^{-C_Ue}\) covering
the \(z_0\)-box, with mesh at most half that radius.  Choose
\(C_U\) using the uniform discriminant derivative bounds and the
lower bounds in \Cref{a02:source-discriminant}.  For each pair,
a net center within half a radius of \(z_0\) then gives a ball
whose fixed complex enlargement \(8U\) has both of that pair's
discriminants nonzero.  Assign the first such ball in a fixed
ordering.  Admissibility is a condition on the pair and the ball;
no ball is required to avoid the zeros of every equation.
The pair's two converted equations have simple holomorphic
branches on \(8U\), real on the real ball.  Choose the nearest
real branch to \(w(t_0)\).
For an unreplaced test \Cref{a02:nearest-quadratic-root}
applies; for a replaced test use its affine graph.
In either case,
\begin{equation}
 |w(t_0)-f(z_0)|\le\delta N^{O(e)}.
 \label{a02:near-source-graph}
\end{equation}

We claim that, apart from power-small pair mass, the two
chosen sheets differ in norm on \(2U\) by at most
\begin{equation}
 \delta N^{C_*e}
 \label{a02:local-sheet-match}
\end{equation}
for a fixed sufficiently large \(C_*\).
To prove this, consider discrepancies with norm between
\(M\) and \(2M\), for dyadic \(M\ge\delta N^{C_*e}\).
Group by \(M,D,U\), world and new block, and by the two
center jet cubes at spacing \(M\), in the scaled coordinates
of \(U\), through a fixed sufficiently high order.
Paired cubes are bounded neighbors by
\Cref{lem:algebraic-norm}.  The same lemma shows that all
members of such a group are within \(O(M)\) of a common
reference sheet in norm.

Each old label has at most \(N^{O(e)}\) group options, with
an exponent independent of \(C_*\).  The \(D\)-compatibility
list is subpower, the fine ball grid has \(N^{O(e)}\) members,
and there are at most two branches per ball and boundedly
many neighboring jet cubes.  The range of \(M\) is polynomial
in \(N\), so it contributes only \(O(\log N)\) choices.
For this last assertion, the chosen branch value near
\(z_0\) is polynomially bounded by the trajectory coefficient
bounds and \Cref{a02:near-source-graph}.  Implicit
differentiation, using the derivative lower bound from
\Cref{a02:source-discriminant}, bounds its fixed jets
polynomially.  Affine replacements use an inverse curvature
only above their cutoff.  Finite-jet norm control bounds
the whole comparison norm.  The polynomial exponent here
may depend on the exact problem; its logarithmic contribution
is absorbed in \(N^e\) for large \(N\).  An unused far
quadratic root is irrelevant to this argument.

Let \(m_G\) be the mass of one group in its world/block.
At a given base point, all its values \(w(t_0)\) lie in a
moving band of width \(O(M)\), using
\Cref{a02:near-source-graph} and choosing \(C_*\) after its
loss constants.  If \(M\le1\), choose a capped grid depth
with \(M\le a_r\le MN^e\).  Only boundedly many hidden bins
of that width meet the band.  The cap-class inequality gives
\[
 n(r)\le n(L)(a_r/\delta)^d.
\]
With \Cref{a02:coarse-jacobian}, this proves
\begin{equation}
 m_G\le N^{O(e)}n_Q(M/\delta)^d.
 \label{a02:matching-group-mass}
\end{equation}
When \(M>1\), use depth zero.  The native value and incidence
derivative tests in version two allow only a subpower number of
unit hidden bins at each exact base, by
\Cref{lem:quadratic-sublevel}; in version one \(w\)
is subpower bounded.  Since \(d\ge0\) for a nonempty cap
class, the depth-zero cap again gives
\Cref{a02:matching-group-mass}.

In each color omit group indices, consisting of a label and
branch, whose marginal mass before the joint deletion is
less than \(n_QN^{-C'e}\).  Choose \(C'\) so that the total
cost, summed using the preceding group-option count and
the old label count, is less than \(N^{-2e}\).
The remaining lists have cardinality
\[
 \le N^{O(e)}(M/\delta)^d
\]
by \Cref{a02:matching-group-mass}.
If the total mass of these large-discrepancy groups were at
least \(N^{-e}\), some group would retain mass
\(m_G'\ge m_G/2\) after both color deletions.
Normalize that group's surviving law to probability.
For a color, start with index weights \(n_Q/m_G'\).
Their sum is at most \(2N^{C'e}\), because each retained
index had its stated pre-deletion floor.  Normalize these
weights to probability.  \Cref{a02:source-base-density}, after rescaling the ball,
then bounds each index-base density by \(N^{O(e)}\) times
its normalized index weight.

The lists now agree in value to \(\delta N^{O(e)}\), have
paired norm discrepancies comparable to \(M\), and satisfy
all other hypotheses of \Cref{a02:local-graph-matching}.
Every displayed loss \(N^{O(e)}\) is a sufficiently small
fixed power of \(M/\delta\) if \(C_*\) is chosen large enough.
The constants determining these losses and the option count
were fixed before \(C_*\).  The corollary gives a
contradiction.  This proves \Cref{a02:local-sheet-match}
outside power-small mass.

\par\medskip\noindent\textbf{2. Representative equations at a common anchor.}

The graphs now agree near their sampled source point, but those
comparison balls vary with the point.  We next record their agreement
at one common complex point for each frame \(D\).  This will give
a bin that the witness label can list and a real representative
equation that fits the source trajectory on the whole new block.
The source time remains averaged throughout this part.

\paragraph{Common complex anchors.}  For each \(D\), including its
world and new-block context, choose a random anchor
\(\alpha_D\) uniformly in a fixed bounded complex box with
interior.  Also choose independently a random
\(\sigma_D\) in a fixed bounded complex disk.
For a source point \(z_0\), use the complex line
\[
 z=z_0+s(\alpha_D-z_0)
\]
and the broken path \(0\longrightarrow\sigma_D
\longrightarrow1\) in its \(s\)-plane.

For each retained pair, with probability
\(1-N^{-\Omega(e)}\) over these choices, both converted
discriminants have lower bounds \(N^{-O(e)}\) on a tube
of radius \(N^{-O(e)}\) around this common path.
Here are the quantitative details.  The coefficient
bounds and the lower bound at \(z_0\) imply a polynomial
norm lower bound on the fixed anchor domain.  Polynomial
sublevel estimates therefore make both values at
\(\alpha_D\) at least \(N^{-B e}\), except with
power-small probability, for a fixed sufficiently large
\(B\).  Restrict a discriminant to the complex line; its
coefficients in \(s\) are bounded by \(N^{O(e)}\).
Its nonzero values at \(0,1\), together with these derivative
bounds, place those endpoints at distance at least
\(N^{-O(e)}\) from all zeros in the bounded path region.

Remove disks of a much smaller radius \(N^{-Be}\), with
\(B\) increased, around the boundedly many such zeros.
A segment from either endpoint to a uniformly chosen
\(\sigma_D\) meets one disk only on a set of power-small
probability: the disk subtends an angle bounded by its
radius divided by its distance from that endpoint, and
the \(\sigma_D\)-density is bounded.  On a path avoiding
these disks, factor the one-variable polynomial and compare
each factor with its value at zero.  Nearby zeros cost
only fixed powers of \(N^{-e}\); far zeros cost fixed
constants.  This gives a lower bound \(N^{-O(e)}\) all
along the path.  The coefficient bounds in the full
\(z\)-variables extend it to the stated smaller tube.
These exceptional probabilities may be averaged against
the original pair law, so their absolute cost is
power-small.

Choose common endpoint comparison balls of radius
\(N^{-O(e)}\), with fixed margins inside these tubes and,
at the source, inside the original stable ball.
Continue the two chosen branches.  First use
\Cref{lem:algebraic-norm} to extend
\Cref{a02:local-sheet-match} to the interior complex source
domain, and then use \Cref{a02:algebraic-path}.
At the anchor their difference, and all required fixed
jets in \(D\)'s \(z\)-coordinates, are at most
\(\delta N^{O(e)}\).

For each \(D\), bin a sufficiently long jet of the source
branch at \(\alpha_D\), separating real and imaginary
parts, at one common spacing
\begin{equation}
 \delta N^{C_b e},
 \label{a02:anchor-spacing}
\end{equation}
where \(C_b\) is fixed after the propagation constants.
On retained pairs these anchor bins are bounded neighbors.
The reached jets have polynomial range: their source
norms have polynomial bounds by the preceding near-branch
argument, and \Cref{a02:algebraic-path} propagates the norm
of the branch itself.  Thus we can prescribe polynomial
ranges for all bins used on successes, even when an
unused far root has much larger coordinates.

Given \(A'\), the source \(D\) belongs to a subpower
compatibility list.  For each such \(D\), enumerate the simple
roots at the anchor of the converted equation
\(\widetilde P_{A',D}\), including its affine replacement when
the early curvature test required one.  There are at most two,
and their jet tuples are determined by those roots.  These are
the same converted sheets used in the local comparison.
Different continuation paths can interchange these two tuples but
cannot create further ones.  Taking the neighboring bins
of both tuples gives a list of only subpowerly many
possible source bins.  This is genuine old-\(p_L\) list
knowledge: its size is subpower, rather than merely
\(N^{O(e)}\).

\paragraph{A representative test for each bin.}
The source label \(A\) now supplies the branch to be fitted on
the new block.  The opposite-time label \(A'\) supplies only the
list just constructed, which makes the eventual output label
visible at \(p_L\).
For each \(D\) and used bin choose a representative
equation \(\widehat P\) among the converted real equations
\(\widetilde P_{A,D}\) from all old source subblocks.  Require the fixed
discriminant coefficient bounds, its prescribed
\(N^{-O(e)}\) lower absolute bound at \(\alpha_D\), and
a branch whose jet lies in that bin.  These are fixed
equation/anchor eligibility tests, allowing affine
replacements.  Choose the first eligible finite option
in a predetermined order.  This table depends on the
anchor, \(D\), and bin, and not on \(t_0\) or the sampled
trajectory.  Every source on a retained success is itself
eligible.

On a common anchor ball with an inverse-power radius,
the representative branch \(\hat f\) is simple and
holomorphic.  For the source branch \(f\) in the same
bin, the finite-jet part of \Cref{lem:algebraic-norm}
gives
\[
 \|f-\hat f\|\le\delta N^{O(e)}.
\]
The derivative of \(\widehat P\) on \(\hat f\) is bounded
by the square root of its discriminant, hence by
\(N^{O(e)}\); also
\(|\widehat P_{ww}|\le\delta^{-1}\).
Taylor expansion in \(w\) therefore yields on an interior
anchor ball
\begin{equation}
 |\widehat P(z,f(z))|\le\delta N^{O(e)},\qquad
 |\widehat P_w(z,f(z))|\le N^{O(e)}.
 \label{a02:anchor-evaluations}
\end{equation}
For example, the quadratic error in the first estimate
is bounded by
\(\delta^{-1}(\delta N^{O(e)})^2
=\delta N^{O(e)}\).

Return these \emph{evaluations} to the source along the
source-safe path.  Each is a fixed-degree polynomial
evaluation of the continued branch of \(\widetilde P_{A,D}\), so
\Cref{a02:algebraic-path} applies.  Only the source
discriminant is required to stay nonzero on this return
path.  In particular, the representative branch is not
continued back, and no assertion is made that it has
a real root near every source point.

On a real segment near \(t_0\) of relative new-block
length at least \(N^{-O(e)}\), the actual line remains in
the source endpoint ball and
\[
 |w-f(z)|\le\delta N^{O(e)}.
\]
For a replaced source this is its affine-center estimate.
For an unreplaced source, \(P_{A,D,w}\) is affine along the
line, has the upper bound in
\Cref{a02:old-test-extrapolation}, and has its own-time
lower bound \(|P_{A,D,w}(z_0,w(t_0))|\ge K_0^{-1}\).
This is the original selected-incidence floor for \(A\);
the opposite-time estimate for \(A'\) is not needed for this return.
Shrink the segment by a fixed power of \(N^{-e}\) so
that this lower bound persists within a constant factor;
shrink it again to fit the source endpoint ball, using
the bounded \(z\)-speed.  The early cutoff \(C_A\)
still makes \(|P_{ww}P|\ll|P_w|^2\) there.
\Cref{a02:nearest-quadratic-root} thus applies
continuously to the same source branch chosen at \(t_0\).
If \(t_0\) is near the edge of the new block, use the
one-sided portion of this segment; it has the same
lower order of length.

Use \(|\widehat P_{ww}|\le\delta^{-1}\) once more to
replace \(f(z)\) by the actual \(w\) in the returned
estimates \Cref{a02:anchor-evaluations}.  Polynomial
extrapolation along the line then gives, on the entire
new block,
\begin{equation}
 \sup_I|\widehat P(z,w)|\le\delta N^{C_2e},
 \qquad
 \sup_I|\widehat P_w(z,w)|\le N^{C_2e}
 \label{a02:representative-upper}
\end{equation}
for a fixed \(C_2\).  In version two all trajectory
coordinates are affine.  In version one the special
form \(w-F(t,y)\) still has degree at most two along a
trajectory.  Thus the extrapolation has fixed degree
in both cases and costs only the indicated fixed powers.

We also need a lower bound for the representative
discriminant at the real source:
\begin{equation}
 |\widehat{\mathcal D}(z_0)|
 =
 |\widehat P_w^2-2\widehat P_{ww}\widehat P|
 \ge N^{-C_3e}.
 \label{a02:representative-disc-lower}
\end{equation}
This fails only on power-small pair mass, with an exponent
\(C_3\) chosen before the later curvature split.  To prove this,
fix the anchors and the resulting representative table.  Each
representative has the prescribed inverse-power discriminant
lower bound at its anchor.  Polynomial norm comparison transfers
it to the fixed real \(z\)-box \(B_{\rm real}\): for an already
fixed exponent \(b_{\rm anc}\), uniformly over the table,
\[
 \|\widehat{\mathcal D}\|_{B_{\rm real}}
 \ge N^{-b_{\rm anc}e}.
\]
Let \(\theta>0\) be a fixed polynomial sublevel exponent for
these degrees and this box.  Then
\[
 \bigl|\{z\in B_{\rm real}:
       |\widehat{\mathcal D}(z)|<N^{-C_3e}\}\bigr|
 \le C N^{-\theta(C_3-b_{\rm anc})e}.
\]

For summation, fix a source label \(A\).  Its \(D\)
is fixed.  At the fixed anchor its converted equation
\(\widetilde P_{A,D}\) has at most two root/jet tuples,
hence at most two source
bins and a bounded number of associated representatives.
The number of possible homotopy classes of a varying
path does not enter this count.  Charge the exceptional
volumes against \Cref{a02:source-base-density} for
this source, then sum using \Cref{a02:old-label-count}.
In each world and block the bad pair mass is at most
\begin{align*}
 &\frac{K_0N^e}{n_Q}
   \bigl(n_QN^{C_1e}\bigr)
   C N^{-\theta(C_3-b_{\rm anc})e}\\
 &\hspace{12mm}\le
 K_0N^{[C_1+1-\theta(C_3-b_{\rm anc})]e}.
\end{align*}
The bounded representative count is absorbed in \(K_0\).
Summing with \(\omega_\gamma q_1\) gives the same global bound.
Choose
\[
 C_3>b_{\rm anc}+\frac{C_1+2}{\theta}.
\]
This makes the exception power-small using only constants
already fixed.  In particular it is independent of the later
curvature threshold.  The charge is made while \(t_0\) is still
averaged, using the representative table's independence of \(t_0\).

\par\medskip\noindent\textbf{3. Construction of the longer atlas.}

The representative equations now satisfy the whole-block bounds
\Cref{a02:representative-upper}.  Outside power-small pair mass,
their discriminants also satisfy
\Cref{a02:representative-disc-lower} at the source.  We use these
facts to construct trajectory assignments on the longer blocks.
The source label will determine an assignment, but will not be
part of its output label.

\paragraph{Fixing the source and assigning anchor bins.}
Let \(\mathbf a\) denote the anchor and detour table, and let
\(\mathcal S_{\gamma,I}(\ell,t_0,t;\mathbf a)\) be the event that
the pair passes all restrictions through
\Cref{a02:representative-disc-lower}.  For
\(\vartheta=(\mathbf a,t_0)\), define
\[
 s_{\gamma,I}(\vartheta)
 =\int\int_I
   \1_{\mathcal S_{\gamma,I}}(\ell,t_0,t;\mathbf a)
       \frac{dt}{q_1}\,d\nu_\gamma(\ell).
\]
Sample \(\vartheta\) using the anchor and detour laws above and
independent uniform \(t_0\in I\).  For each world and block,
choose \(\vartheta_{\gamma,I}=(\mathbf a_{\gamma,I},t^{\rm src}_{\gamma,I})\)
with mass at least its average.  The preceding estimates give
\[
 M_{\rm fr}:=
 \sum_{\gamma,I}\omega_\gamma q_1
      s_{\gamma,I}(\vartheta_{\gamma,I})
 \ge
 \sum_{\gamma,I}\omega_\gamma q_1
      \E_\vartheta s_{\gamma,I}(\vartheta)
 =N^{-o(1)}.
\]
The lower bound concerns this weighted sum; individual blocks may
have much smaller mass.  In a fixed world and block, write
\(t_0=t^{\rm src}_{\gamma,I}\).  The frozen successful measure is
\[
 d\mu^{\rm fr}_{\gamma,I}(\ell,t)
 =\1_{\mathcal S_{\gamma,I}}
   (\ell,t_0,t;\mathbf a_{\gamma,I})
       \,d\nu_\gamma(\ell)\,\frac{dt}{q_1}.
\]
It remains an unnormalized restriction of the original product
prior.  We now define geometric assignments containing its
successful trajectories.

At the fixed source time, require old selected source-chart
membership.  This determines the old label \(A\), its coarse
frame \(D\), and the nearest simple branch of
\(\widetilde P_{A,D}\).  Require its discriminant to be nonzero
on the fixed broken path,
with the prescribed simple-branch conditions at its endpoints,
and continue the source branch to the anchor.  Require its value
and jets to lie in the stored polynomial ranges.  Let \(b\) be
its anchor-jet bin, and require that this bin have a representative
table entry.  Retain only trajectories satisfying the whole-block
spatial bound, both bounds in \Cref{a02:representative-upper}, and
\Cref{a02:representative-disc-lower} at \(t_0\), for the table
entry \(\widehat P_{D,b}\).  Every frozen success satisfies these
conditions.

Assign each such trajectory the preliminary label
\[
 c_0=(D,b),
\]
merging all source labels that give the same pair.  The fixed
source time and disjoint old source assignments give at most one
label per trajectory in each new block.  In particular this
assignment is independent of the later time \(t\).  Its
representative equation is the fixed table entry
\(\widehat P_{D,b}\); the old source identity is no longer
recorded.

For these assignments retain the original later incidences whose
old terminal label \(A'\) supplies \(D\) and \(b\) in its
compatibility lists.  This selection contains the frozen
successes, has weighted mass at least \(M_{\rm fr}\), and is still
a restriction of \(d\nu_\gamma\,dt/q_1\).  It is specified by the
old incidence predicates, the fixed tables, and geometric tests;
the analytical density masks are not part of the selection.
At \(p_L\), the old label \(A'\) belongs to a subpower list.
Each such label gives a subpower list of frames, and for each frame
at most two anchor-jet tuples with boundedly many neighboring bins;
intersect this list with the stored polynomial bin range.
Thus the preliminary label \(c_0\) is known by \(p_L\).

It remains to turn each representative equation into an admissible
real test.  The two curvature cases below preserve these
assignments; only the low-curvature case refines their labels.

\paragraph{Normalizing the representative tests.}
The bounds \Cref{a02:representative-upper} and the choice of
\(C_3\) are already fixed.  Choose \(C_4>C_3/2+2\), then choose
\[
 C_s>C_2+C_4+1,
\]
and finally choose
\begin{equation}
 T>\max(C_2+C_3+1,\ C_s+C_4+1).
 \label{a02:late-constant-order}
\end{equation}
First consider a representative with
\[
 |\widehat P_{ww}|\le\delta^{-1}N^{-Te}.
\]
On every assigned trajectory, the product of curvature and value
in \Cref{a02:representative-upper} is at most
\(N^{(C_2-T)e}\), smaller than the discriminant lower bound in
\Cref{a02:representative-disc-lower}.  At the real source,
\[
 \widehat P_w^2
 =\widehat{\mathcal D}+2\widehat P_{ww}\widehat P.
\]
It follows that the discriminant is positive and
\(|\widehat P_w(z_0,w(t_0))|\ge N^{-C_4e}\).
Consequently the trajectory's block maximum satisfies
\[
 N^{-C_4e}\le M_\partial:=\sup_I|\widehat P_w(z,w)|
 \le N^{C_2e}.
\]
Let \(v\) be the lower endpoint of the dyadic bin containing
\(M_\partial\), refine the label to \(c=(D,b,v)\), and use the test
\(\widehat P_{D,b}/v\).  These refinements remain disjoint
trajectory assignments throughout the block.  The full range of
\(v\) contains only \(O(\log N)\) values, so the refined label
is still known by \(p_L\).  In version one the derivative is
identically one, and we take \(v=1\).

The normalized derivative has block maximum between fixed positive
constants.  The value bound is at most
\(\delta N^{(C_2+C_4)e}\), and the curvature bound is
\[
 \delta^{-1}N^{(-T+C_4)e}.
\]
By \Cref{a02:late-constant-order}, these fit the thickness
\[
 \delta_s=\delta N^{C_s e}.
\]
The derivative upper bound is of order one.  For its incidence
lower bound, omit later times where its absolute value is below
an inverse-subpower number \(\kappa\).  On each assigned
trajectory this derivative is affine and has order-one block
maximum, so the omitted set has relative length \(O(\kappa)\).
Disjointness of the assignments gives total weighted loss at most
\(O(\kappa)\) against the original trajectory and time priors.
The same bound holds for the selected submeasure.  Choose
\(\kappa=N^{-o(1)}\) with \(\kappa=o(M_{\rm fr})\); the loss is
then negligible compared with the available mass.  This estimate
uses no density bound for a law conditioned on pair success.

For the remaining representatives,
\[
 |\widehat P_{ww}|>\delta^{-1}N^{-Te},
\]
use the real affine critical center and the test
\(w-F_{\rm ctr}(z)\), keeping the label \(c=c_0\).
By \Cref{a02:critical-center,a02:representative-upper},
\[
 |w-F_{\rm ctr}(z)|\le\delta N^{(T+C_2)e}
\]
on the entire new block.  This fits thickness
\(\delta_h=\delta N^{(T+C_2+1)e}\), with derivative identically
one and zero curvature.  A real representative root near the
source is not needed in this case.

Choose a curvature case carrying at least half the remaining
mass, and pass to a subsequence if necessary to obtain a common
thickness exponent.  All constants were chosen in the order
\(C_A\), the stable-ball and comparison constants, the continuation
and binning constants giving \(C_2\), then \(C_3,C_4,C_s,T\).
In particular \(C_3\) is independent of \(T\): its sublevel
estimate was proved for the representative table while the source
time was still averaged.

Let \(C\) exceed both final thickness exponents and all needed
smallness constants.  Restrict \(e<cL\), with fixed \(c>0\) small
enough that \(L-Ce>0\) and \(Re\le\eta L\).  The chosen thickness
is \(a_{L'}\), with \(0<L'<L\) and \(L-L'\le Ce\).
The block length \(q_1\) has exponent \(h-e\).

\paragraph{Finite encoding and return to the original coordinates.}
The assignments just constructed use polynomially many old labels,
coarse boxes, anchor bins, and representative table entries.  The
chosen source branches and their propagated jets have polynomial
ranges, as established above; unused far roots need not be stored.
The continuation predicates have bounded semialgebraic complexity
by \Cref{a02:algebraic-path}, and fixed-order implicit
differentiation gives the same property for the endpoint jets.
Nearest-root choices, geometric inequalities, and affine block
maxima also have fixed-degree descriptions.  A whole-block bound
quantifies only over one time variable for each individual test.
Combining these conditions with the old predicates and finite
tables therefore preserves temperedness.

All equation and coordinate coefficients have polynomial bounds.
The affine-center constructions invert curvatures only above their
specified cutoffs, and the low-curvature normalization divides only
by numbers at least a constant times \(N^{-C_4e}\).
Thus the final tests and inverse coordinate changes remain
polynomially controlled.  The measurable choices made during the
proof enter only through the fixed source times, anchors, equations,
and tables, not through analytical masks or their searches.

Finally compose the tests with the real affine coordinate map
\(z=z(t,y)\) and undo \(w\mapsto Bw\).  In version one this
preserves the form \(w-F(t,y)\); in version two it preserves total
degree at most two and restores the derivative unit \(B\).
The moving spatial frames contain the assigned full traces and
have largest singular value at most \(q_1\) times a subpower
factor.  We have constructed tempered disjoint trajectory
assignments with the required whole-block fits, inverse-subpower
selected incidence mass, and labels known by the old \(p_L\).
They form the asserted known atlas.
\end{proof}
